\begin{table}[t]
  \centering
  \small
  \begin{tabular}{lrrrcrr}
    \toprule
    Qué & Imág. & Media & Mediana & Sin localizar & Chicos & Grandes \\
    \midrule
    capa 4 (\texttt{C3k2}) & 30 & 1.59 & 1.09 & 0/30 & --- & --- \\
    capa 6 (\texttt{C3k2}) & 30 & 1.38 & 1.01 & 0/30 & --- & --- \\
    capa 8 (\texttt{C3k2}) & 30 & 1.32 & 0.66 & 9/30 & --- & --- \\
    capa 9 (\texttt{SPPF}) & 30 & 1.56 & 0.75 & 3/30 & --- & --- \\
    capa 10 (\texttt{C2PSA}) & 30 & 0.95 & 0.79 & 5/30 & --- & --- \\
    \texttt{cis\_test} & 120 & 1.46 & 1.12 & 52/120 & 1.89 & 1.43 \\
    \texttt{trans\_test} & 120 & 1.25 & 1.00 & 62/120 & 0.89 & 1.27 \\
    \bottomrule
  \end{tabular}
  \caption{Razón de masa de Grad-CAM: fracción de la activación que cae dentro de las cajas anotadas, dividida por el área que esas cajas ocupan; vale 1.0 si el mapa es uniforme, es decir si el modelo no localiza nada. Arriba, el barrido que decidió qué capa del backbone se usa (`sin localizar' cuenta los mapas idénticamente nulos dentro del animal). Abajo, la razón por partición con la capa elegida (`sin localizar' cuenta las imágenes con razón $\leq 1.0$). `Chicos' son los animales que ocupan menos del 1\% del cuadro.}
  \label{tab:gradcam}
\end{table}
